This chapter presents the artefacts the empirical evaluation rests on:
the reference memory bus, the compressor framework, the multi-fragment
\emph{C1} coordination benchmark, the coordination-success and
critical-token-recall metrics that the cliff analysis turns on, the
precomputed compression cache that decouples compression from
evaluation, and the single-command reproducibility envelope. It gives
a reader enough detail to reproduce the system from the reference
repository without reading the code, and makes the design choices
explicit so the results of \charef{experiments} can be read against
the constraints those choices imposed.

A scope reminder bounding every claim that follows: the planner is a
deterministic regex information-extraction solver for the cliff sweep
(H1, H2) and a single LLM call with all compressed fragments visible
for the model-scaling sweep (Corollary~1, frontier, real-data
transfer). The memory bus is \emph{designed for} multi-agent
integration via the AutoGen back-end in the source tree, but
multi-agent operation is outside the empirical scope
(Section~\ref{sec:scope}).

\section{The Memory-Bus Reference Implementation}
\label{sec:architecture}

\subsection{Three-layer architecture}

The reference implementation organises the memory bus into three
layers that match the access pattern of a fragment-passing workflow:

\begin{enumerate}
\setlength\itemsep{2pt}
\item An \textbf{access layer}, exposing \texttt{read},
\texttt{write}, \texttt{subscribe}, and \texttt{audit} endpoints via
FastAPI and enforcing per-request access-control checks through a
policy middleware.
\item A \textbf{compression layer}, providing a pluggable
\texttt{Compressor} protocol behind which six concrete variants live:
an \emph{identity} no-compression control, the
\emph{LLMLingua-2}~\cite{pan2024llmlingua2} token classifier, an
\emph{instruction-aware filter} based on TF--IDF pruning plus a
cross-encoder re-ranker, a \emph{Phi-3-Mini extractive} span selector
implemented over a local Ollama endpoint, a \emph{truncation}
prefix-keep baseline, and \emph{CAAC}, a cliff-aware adaptive
wrapper that backs off compression when per-fragment recall would
violate the compounding-error model's safety threshold
(Section~\ref{sec:compressors_caac}, full discussion in
\charef{summary}).
\item A \textbf{storage layer}, providing three protocols
(\texttt{AuditLog}, \texttt{Scratchpad}, and \texttt{VectorStore})
implemented for the evaluation by SQLite in WAL mode with an
append-only audit table, an in-memory time-to-live dictionary, and
FAISS-CPU with an HNSW index.
\end{enumerate}

Figure~\ref{fig:architecture} summarises the three layers. Each
component dependency in the source tree is one-directional: the
compression layer does not import from the access layer, and the
storage layer does not import from either. This layering means the
storage backends can be replaced without touching the compression or
access layers.

\begin{figure}[htbp]
\centering
\begin{tikzpicture}[
    layer/.style={draw, rounded corners=2pt, minimum width=12cm,
                  minimum height=1.6cm, align=center, font=\small},
    cmp/.style={draw, rounded corners=1pt, minimum width=1.9cm,
                minimum height=0.6cm, font=\footnotesize},
    arrow/.style={-{Latex[length=2mm]}, thick},
    label/.style={font=\footnotesize\itshape}
]
% Access layer
\node[layer, fill=black!4] (access) at (0,4)
  {\textbf{Access layer (FastAPI)} \\[1pt]
   \footnotesize POST /v1/write \quad GET /v1/read \quad POST /v1/subscribe
   \quad GET /v1/audit \\
   \footnotesize PolicyMiddleware \quad StructLogMiddleware \quad OtelMiddleware};

% Compression layer
\node[layer, fill=black!4, minimum height=2.0cm] (compress) at (0,1.8)
  {\textbf{Compression layer} \\[2pt]
   \texttt{Compressor} protocol \\[2pt]};
\node[cmp] (id)     at (-5.0, 1.5) {identity};
\node[cmp] (lingua) at (-2.9, 1.5) {LLMLingua-2};
\node[cmp] (filter) at (-0.8, 1.5) {filter};
\node[cmp] (phi)    at ( 1.3, 1.5) {Phi-3 ext.};
\node[cmp] (trunc)  at ( 3.4, 1.5) {truncation};
\node[cmp] (caac)   at ( 5.5, 1.5) {CAAC (wrap)};

% Storage layer
\node[layer, fill=black!4, minimum height=2.0cm] (storage) at (0,-1.0)
  {\textbf{Storage layer} \\[2pt]};
\node[cmp, minimum width=3cm] (audit) at (-3.7, -1.3) {SQLite AuditLog};
\node[cmp, minimum width=3cm] (scr)   at ( 0.0, -1.3) {InMemory Scratchpad};
\node[cmp, minimum width=3cm] (vs)    at ( 3.7, -1.3) {FAISS VectorStore};

% Arrows between layers
\draw[arrow] (access) -- (compress);
\draw[arrow] (compress) -- (storage);

% Side annotations
\node[label] at (7.6,  4.0) {ACL-checked};
\node[label] at (7.6,  1.8) {pluggable};
\node[label] at (7.6, -1.0) {append-only};
\end{tikzpicture}
\caption{Three-layer architecture of the memory-bus reference
implementation. Each layer depends only on the layer below it, and
each storage component is hidden behind a \texttt{Protocol} so that
the storage backends can be replaced without touching the access or
compression layers. CAAC is drawn inside the compression layer for
convenience but is a wrapper: it composes any of the other
compressors with a recall-side back-off rule
(Section~\ref{sec:compressors_caac}).}
\label{fig:architecture}
\end{figure}

The three-layer structure is a direct extension of
MemIndex~\cite{saleh2025memindex}, the FCG group's baseline architecture
for distributed agent memory. The two material differences are that
compression is promoted to a first-class layer rather than treated as
a downstream optimisation, and that per-slot tags are part of the
data model rather than out-of-band metadata.

\subsection{Data flow for a write}

A \texttt{POST /v1/write} carries the fragment text, its tag vector,
and an optional task hint. The policy middleware parses the principal
from the \texttt{X-M6-Principal} header and checks that its
access-control mask is a superset of the fragment tag's and its
classification at least the fragment's. On success the service
compresses the fragment to a \texttt{CompressedSlot}, stores it in
the scratchpad, adds any embedding to the FAISS index, and appends a
\texttt{WRITE}/\texttt{OK} audit row whose
$\texttt{chain\_hash}=\textsc{SHA-256}(\texttt{prev\_hash}\mathbin\Vert\texttt{payload\_hash})$;
the response returns the slot id, hex chain hash, and audit row id. A
failed check appends a \texttt{DENY} row (slot \texttt{deny-<uuid8>})
with a $60$-second per-(subject, fragment-id) deduplication so a
client hammering denied writes cannot inflate the log.

\subsection{Data flow for a read}

A \texttt{GET /v1/read/\{slot\_id\}} flows symmetrically: the
middleware attaches the principal, the service looks up the slot,
checks the principal--tag relationship, and either returns the slot
with its audit row id or raises a policy-denied / slot-not-found
error. Read misses are audited as \texttt{READ}/\texttt{ERROR} rows
with the same $60$-second per-(subject, slot-id) deduplication.
Appendix~\ref{appendix} gives an end-to-end \texttt{curl} trace.

\subsection{Tags, classifications, and the audit log}

Every fragment and every compressed slot carries a tag vector
recording its provenance and access-control state. The tag vector has
four fields: a \texttt{uint64} access-control mask, a five-tier
classification level (\textsc{public} $<$ \textsc{internal} $<$
\textsc{confidential} $<$ \textsc{restricted} $<$ \textsc{secret})
drawn from enterprise data-classification practice (e.g., ISO/IEC 27001
organisational schemes); the broader risk-management framing follows
the NIST AI Risk Management
Framework~\cite{nistairmf}, but the five-tier lattice itself is
enterprise convention and not prescribed by NIST AI RMF (which
defines Govern/Map/Measure/Manage functions rather than a
labelled classification lattice). The remaining two fields are a
tuple of source identifiers recording provenance and a tuple of
parent-slot identifiers recording which previously-compressed slots
this slot inherits from. When slots are
merged, the union's access-control mask is the bitwise OR of the
parents' masks and the union's classification is the maximum; the
policy predicate is symmetric for reads and writes.

The audit log is an append-only SQLite table; each row records the
\texttt{event\_type} (WRITE/READ/SUBSCRIBE/DENY/COMPRESS/EVICT),
\texttt{slot\_id}, requester ACL mask and classification, result, and
a tamper-evident hash chain
($\texttt{chain\_hash}=\textsc{SHA-256}(\texttt{prev\_hash}\mathbin\Vert\texttt{payload\_hash})$).
An \texttt{INSTEAD OF UPDATE} and an \texttt{INSTEAD OF DELETE}
trigger refuse any mutation of an existing row, and a unit test drops
the trigger, edits a \texttt{payload\_hash} at the byte level, and
confirms that the \texttt{verify()} chain-walker reports failure.

A second SQLite table records every paid model call (provider,
model, token counts, EUR cost, wallclock), sharing the audit database
for transactional consistency. Pricing follows a representative
frontier rate (Anthropic Claude Sonnet tier, USD~$3$/$15$ per million
input/output tokens, $\approx$
EUR~$2.76$/$13.80$ at $0.92$~USD/EUR) against an on-premises marginal
$\approx$ EUR~$0.05$ per million tokens~\cite{fcgfinancial2026}; the
ledger feeds the EUR-per-workflow cost model of the RAG catalogue
(Section~\ref{sec:rag_pipelines}).

\section{Retrieval-Augmented Generation Pipeline Catalogue}
\label{sec:rag_pipelines}

The H3 analysis (Appendix~\ref{appendix:rag}) compares three
pipeline placements over the shared FAISS-CPU + HNSW +
\texttt{bge-large-en-v1.5} stack: \textbf{P1} compress-first
(compress the corpus, then index it), \textbf{P2} retrieve-first
(index the full corpus, then compress the top-$k$, the classical
LongLLMLingua~\cite{jiang2024longllmlingua} configuration), and
\textbf{P3} a joint relevance router (pass high-scoring fragments
verbatim, compress mid-scoring, discard low-scoring; after
ACC-RAG~\cite{guo2025dynamic}). Because the pre-specified sign-flip
did not appear and the cost model does not meter compression compute,
H3 is reported as a secondary result in Appendix~\ref{appendix:rag}.

\section{Compressor Framework}
\label{sec:compressors}

Every compressor variant conforms to a single Python \texttt{Protocol}
with four methods: \texttt{compress(fragment, target\_ratio)}
returning a \texttt{CompressedSlot}; \texttt{decompress(slot)}
returning a best-effort text reconstruction; \texttt{embed(slot)}
returning an optional embedding for vector-store indexing; and
\texttt{model\_card()} returning a structured card containing the
compressor identifier, family, base model, default target ratio, and
provenance hashes. The thesis ships six concrete implementations.
The identity compressor is the no-compression control that every
experiment needs to guard against runtime-bug confounds.

The remaining five compressors fall into four families: a
token-level hard-prompt filter (LLMLingua-2); an instruction-aware
heuristic filter; an extractive small-LLM span selector
(Phi-3-Mini extractive); a trivial prefix-truncation baseline; and a
cliff-aware adaptive wrapper (CAAC). \emph{All five are
training-free}, which is a deliberate design choice: the thesis
isolates the compression mechanism from any training-data effect
and lets the cross-compressor comparison turn on the algorithm
rather than on training-budget asymmetry.

\subsection{LLMLingua-2 (token-level classifier)}

The LLMLingua-2 wrapper is a thin adapter around the
\texttt{llmlingua} package~\cite{pan2024llmlingua2}. Construction
takes a target compression ratio; the package is asked to retain
$1/r$ of the tokens, scored by a small XLM-RoBERTa classifier
distilled from a more capable teacher. The force-tokens list
(sentence boundary marks) is validated against the wrapped
tokenizer's vocabulary at initialisation and unknown tokens are
dropped with an explicit log line. PyTorch runs the wrapped
XLM-RoBERTa model on both CPU and Apple's Metal Performance Shaders
backend, so the compressor runs natively on Apple Silicon without
modification.

The empirical token-recall curve $q(r)$ produced by LLMLingua-2
under critical-token-recall measurement (Section~\ref{sec:ctr})
is smooth and monotonic across the sweep ratios used in
\charef{experiments}, which is the property the
compounding-error model presumes.

\subsection{Instruction-aware filter}

The instruction-aware filter is a pure-Python heuristic baseline.
It operates in two stages. The first is a sentence-level filter
that ranks sentences by their relevance to the task hint using a
cross-encoder reranker (\texttt{BAAI/bge-reranker-base}) when
available and a lexical-overlap fallback when not. The second is a
token-level trim that drops English stop-words and short tokens
($\le 2$ characters) until the target compression ratio is reached.
The stop-word list is encoded explicitly so that the trim is
deterministic given the task hint, the fragment, and the target
ratio: a subtle bug showed that dropping stop-words in Python
set-iteration order introduced run-to-run variance.

The filter's reranker preserves task-keyword tokens that simpler
token-level methods drop, which is why its $q(r)$ curve on the
multi-step retrieval family (Section~\ref{sec:c1}, family-c) decays
gracefully rather than collapsing: chain-reference tokens like
``entry'' and ``FINAL'' survive aggressive compression.

\subsection{Phi-3-Mini extractive span selector}
\label{sec:phi3}

The Phi-3-Mini extractive compressor is the only thesis compressor
that uses an LLM at inference time. It runs the Phi-3-Mini-3.8B
instruction-tuned model~\cite{phi3} via a local Ollama
endpoint~\cite{ollama} with a structured prompt that asks for
verbatim contiguous spans from the source. The post-hoc verifier
\texttt{Phi3ExtractiveCompressor.\_verify\_extractive} checks the
output: every token in the output must appear in the source. Two
loosenings of the strict verifier were necessary in practice. First,
a $15\%$ novel-token tolerance: small language models frequently
insert function words (``the'', ``is'', ``and'') even when
explicitly asked not to. Second, a post-hoc
\texttt{\_strip\_novel\_tokens} pass that removes any token not in
the source after the verifier passes; this guarantees zero novel
tokens in the final output even if the model produced some during
generation. Outputs that fail even after stripping fall back to
LLMLingua-2 at the same target ratio.

Phi-3-Mini extractive has a \emph{compression ceiling}: span-level
selection, the stripping pass, and the LLMLingua-2 fallback together
bound the achievable ratio at about $2.5\times$, regardless of the
requested target. This is documented as a limitation in
\charef{summary} and is the reason all evaluation reports
\emph{achieved} compression ratio (the ratio actually delivered)
alongside the requested target ratio.

\subsection{Truncation (prefix-keep baseline)}

The truncation compressor keeps the first $1/r$ fraction of tokens
by whitespace split. It is deterministic by construction and
achieves the requested ratio exactly. Its purpose is the
lower-bound argument: every more sophisticated compressor must
beat truncation on the coordination metric to justify its cost.

Truncation has a peculiar property: its single-agent
question-answering F1 is competitive with the more sophisticated
compressors at low to moderate compression ratios because the
prefix tokens are naturally representative. Its coordination
success collapses earlier than the others, because the tokens
removed at the tail include the task-critical content that the
planner needs to combine. This dissociation is part of the H1
finding: a token-overlap quality metric over-reports the utility
of even the most naive compressor.

\subsection{Cliff-Aware Adaptive Compression (CAAC)}
\label{sec:compressors_caac}

CAAC is a wrapper, not a fifth compressor family. Given an inner
compressor (LLMLingua-2 in the headline experiments; any of the
others in the ablation) and a per-family recall threshold
$\theta_q$, CAAC runs the inner compressor at the requested target
ratio, measures the resulting critical-token-recall, and if recall
falls below $\theta_q^{1/N}$ binary-searches downward for the
largest ratio that keeps recall above the threshold. A
\texttt{min\_ratio} floor of $1.5\times$ prevents CAAC from
abdicating to no-compression on hard fragments. The full
algorithmic and empirical discussion lives in
\charef{summary} together with the constructive framing
that CAAC is the operational realisation of the compounding-error
bound. Here it suffices to note that CAAC consumes the same
\texttt{Compressor} protocol as the other variants, which is what
makes its drop-in evaluation in \charef{experiments}
straightforward.

\section{The C1 Coordination Benchmark}
\label{sec:c1}

\subsection{Why a synthetic benchmark}

The coordination cliff $\tau^{*}$ is a position on a curve, and
detecting it requires sweeping ratios densely from $1\times$ to
$16\times$ with everything else held constant. Real benchmarks
(HotpotQA, MultiHop-RAG) confound the cliff with task difficulty,
retrieval distribution, and per-instance variance, so the thesis
characterises the cliff on a synthetic generator first (where
information density is placed by construction) and validates the
transfer on real benchmarks (Corollary~2). The C1 generator produces
$150$ instances across three families from a single seed, with the
config hash recorded in the manifest so any output CSV traces back to
its generator inputs.

\subsection{Three workload families}
\label{sec:c1_families}

\begin{description}\setlength\itemsep{2pt}

\item[Family-a: cross-document fact aggregation.] Eight source
documents, one per Vignette~3.7 institutional
system~\cite{fcgusecase2026}, each carrying a recorded-hours and an
approved-budget number from a seed-fixed distribution; the planner
aggregates both across the eight, scored by a regex extraction
pipeline against the deterministic arithmetic sum. Family-a is the
densest ($\theta_\text{info}\approx0.97$): every multi-digit number
is task-critical and the cliff is sharp, so it is where the
compounding-error model is most directly testable.

\item[Family-b: constraint-satisfaction planning.] $N$ sub-tasks
with per-task loads and $K$ capacity-bounded workers; the planner
must assign every sub-task within capacity. The generator guarantees
at least one feasible assignment (capacity-respecting greedy), and
the scorer is a \emph{feasibility checker}: any valid assignment
scores~$1$. Constraint tracking is hard for $\le8$B planners even
without compression, so the family-b baseline $p_0$ is low for small
planners and the cliff is undetectable there: a \emph{floor effect}
\charef{experiments} treats explicitly.

\item[Family-c: multi-step retrieval.] A linear chain of fragments,
each carrying a pointer to the next (\emph{entry~X}) or, at the leaf,
the answer (\emph{FINAL-XXXX}); chain lengths $2$--$4$, padded with
$4$--$8$ distractor sentences, scored against the leaf answer.
Family-c is the most distributed of the three: many tokens are distractors,
information density $\theta_\text{info}\approx 0.62$ (distinct from the recall
threshold $\theta_q$; see Section~\ref{sec:theta_info_def}), and the
cliff is gradual; it is the family on which the instruction-aware
filter's reranker, which preserves chain-reference keywords, can
avoid a detectable cliff entirely.
\end{description}

Each family contributes $50$ instances. The three families together
span a deliberate range of information densities so that the cliff
mechanism (a recall threshold $\theta_q$) and the cliff-position
shift mechanism (information density $\theta_\text{info}$) can be
separated; this separation is what \charef{experiments}
formalises as Corollary~2.

\subsection{Tag distributions for the privacy axis}

The C1 generator supports three synthetic tag distributions
(\emph{uniform}, \emph{skewed}, \emph{hierarchical}) used by H4 to
probe the governance-stringency dimension. The hierarchical
distribution sets higher classification levels to imply strict
supersets of the lower-level bits, matching enterprise tag
hierarchies and letting H4 disentangle classification level from
access-control mask.

\subsection{Coordination success and supporting metrics}

The coordination scorer is a pure function of the workload and the
planner's output; it does not call any LLM. It returns one primary
metric (a per-instance binary \emph{coordination success}) and
four supporting metrics: the F1 over extracted answer tokens versus
ground truth, the achieved compression ratio
($|\text{source}|/|\text{compressed}|$), generic token-recall (the
fraction of source tokens preserved in the compressed text), and
\emph{critical-token-recall} (CTR; Section~\ref{sec:ctr}).
Aggregation across workloads $\times$ seeds is by arithmetic mean;
$95\%$ confidence intervals are computed by workload-level
bootstrap, never by per-row bootstrap which double-counts the seed
dimension.

The reason the scoring pipeline never calls an LLM is the same
reason H1 and H2 use a deterministic regex parser as the planner:
isolating the compression-quality signal from LLM run-to-run
variance. The trade-off is that the cliff measured here is an upper
bound on what an LLM-based scorer would measure; the trade is in
favour of clean attribution rather than ecological validity, and
\charef{summary} discusses the limitation.

\subsection{The critical-token-recall metric}
\label{sec:ctr}

A standard token-recall measurement
$q_\text{token}(r) = |T_\text{compressed} \cap T_\text{source}| /
|T_\text{source}|$ counts every preserved token equally. This proved
a poor proxy for task-relevant information preservation: the
Phi-3-Mini extractive compressor at low ratios preserves common
function words like ``the'' and ``and'' but drops some critical
numeric tokens. Its token-recall is high; its coordination success is low.

The critical-token-recall (CTR) metric is the family-specific
restriction:
\[
  q_\text{CTR}(r) = \frac{|T_\text{compressed}^{\,\text{crit}} \cap
                          T_\text{source}^{\,\text{crit}}|}
                         {|T_\text{source}^{\,\text{crit}}|},
\]
where $T^\text{crit}$ is the set of \emph{task-critical} tokens for
the family:
\begin{itemize}\setlength\itemsep{1pt}
\item Family-a: multi-digit numeric tokens (at least two digits, to
exclude single-digit noise);
\item Family-b: all numeric tokens (load and capacity values);
\item Family-c: chain-reference tokens (\emph{entry-X} patterns) and
the literal \emph{FINAL} marker.
\end{itemize}

CTR is what the compounding-error model uses as the per-round
retention rate $q(r)$; CAAC uses CTR as its back-off signal; and
the predicted-versus-empirical $\tau^{*}$ figure in
\charef{experiments} uses CTR to derive $\theta_q$. Generic
token-recall is reported for backwards compatibility with prior
analyses but is not used in the verdict pipeline.

\section{Compression Cache}
\label{sec:cache}

Compression is the slowest evaluation step (Phi-3-Mini extractive
is a $\sim11$~second Ollama call per fragment), and the cliff sweep
touches $\approx30{,}000$ (compressor, ratio, workload, task-hint)
cells, so a precomputed cache (\texttt{CompressionCache} /
\texttt{CachedCompressor} in \texttt{src/m6/compressors/cache.py},
keyed on
$(\text{compressor}, r, \text{fragment\_id}, \text{hash}(\text{task\_hint}))$)
is populated once on the GPU and consumed as JSON thereafter. This
decouples compression from evaluation, letting H1, H2, the frontier
sweep, and the CAAC ablation run on a laptop, and guarantees every
run sees identical compressed text for a cell, so run-to-run
differences are protocol, not compressor stochasticity (cache hit
rate $100\%$ by construction).

\section{Inference Backends}
\label{sec:inference}

The unified \texttt{InferenceBackend} protocol exposes a single
\texttt{complete(prompt, max\_tokens, temperature, stop)} method
together with an optional \texttt{embed(text)} method. Three concrete
backends are used in the evaluation. The first is a local Ollama
endpoint~\cite{ollama} that hosts Qwen-2.5-Instruct at the three
local sizes (\textsc{1.5B}, \textsc{3.8B}, \textsc{8B}) used in
the model-scaling sweep, Phi-3-Mini-3.8B for the extractive
compressor, and Llama-3.1-8B-Instruct~\cite{llama31modelcard} as the
local protected-fact reader. The second is the OpenAI-compatible Featherless API used
for the frontier validation (Qwen-2.5-72B-Instruct, DeepSeek V4
Pro). The third is the OpenAI API for the GPT-class frontier
arm (excluded from the calibrated regime per
\charef{summary}, retained on disk as a noncanonical
diagnostic). All API calls are routed through the cost ledger so
that the EUR-per-workflow figures in
\charef{experiments} are auditable.

\section{Design Choices and Alternatives Considered}
\label{sec:design_choices}

The artefacts above are not a default stack; each tool and parameter
was a choice with a credible alternative, catalogued here so the
results of \charef{experiments} can be read against the decisions
that made them. The per-hypothesis falsifiability bars are discussed
at the point of pre-specification
(\charef{experiments}~Section~\ref{sec:prereg_refinement}).

\subsection{Tooling: storage, retrieval, and serving stack}
\label{sec:design_tooling}

The serving stack is chosen for single-host reproducibility:
\textbf{FAISS-CPU}~\cite{faiss} with an HNSW index (over Qdrant/Chroma/pgvector)
so the reference service ships as one Docker image with no external
database and no GPU dependency for retrieval;
\texttt{BAAI/bge-large-en-v1.5} embeddings with a
\texttt{bge-reranker-base} re-ranker (open-weight, digests in
\texttt{docs/MODEL\_CARDS.md}); \textbf{FastAPI}/Starlette for the
access layer; \textbf{SQLite-WAL} for the audit log (one-file
footprint, concurrent reader--writer, a replicated multi-host chain
named as future work); \textbf{Ollama} for local LLM serving
(uniform model lifecycle across the Apple-Metal and CUDA hosts behind
one \texttt{complete()} call, Section~\ref{sec:inference}); and the
OpenAI-compatible \textbf{Featherless} endpoint for the frontier arm,
chosen for its flat per-token price so the EUR cost model compiles
without per-provider markup tracking.

\subsection{Compressors considered but not swept}
\label{sec:design_compressors_declined}

Three compressor families from the survey were considered and
excluded, recorded so the omissions are deliberate.
\textbf{Selective Context}~\cite{li2023selective} requires a scoring
forward pass at compression time (asymmetric compute with the swept
four) and its design point (task-aware lexical pruning with no
end-to-end LM) is already covered by the instruction-aware filter.
\textbf{RECOMP}~\cite{xu2024recomp} and \textbf{learned compression}
(gist tokens~\cite{mu2023gist},
AutoCompressor~\cite{chevalier2023autocompressor},
ICAE~\cite{ge2024icae}) require fine-tuning, which would confound the
deliberately training-free cliff comparison with training-budget
asymmetry (re-evaluating them is named as follow-on work).
\textbf{xRAG}~\cite{cheng2024xrag} represents a document as a single
token ($>100\times$), beyond the $1$--$16\times$ regime where the
cliff sits. For the extractive compressor's base model, Phi-3-Mini
was chosen over Llama-3.2-3B and Mistral-7B on pilot
extractive-verifier pass rates ($\sim 75\%$ vs
$\sim 62\%$/$\sim 71\%$ pre-fallback) at the smallest deployment
footprint.

\subsection{Threshold and sweep-grid parameters}
\label{sec:design_thresholds}

The cliff measurement turns on a few numerical parameters. The
\textbf{ratio grid} $\{1,2,3,4,5,6,8,10,12,16\}\times$ is denser at
the low end where the dense-family cliff sits; a fine-grid
re-resolution (\charef{experiments}~Section~\ref{sec:h2}) later
showed the $r\in\{1,2\}$ gap is itself too coarse for the
deterministic-solver family-a cliff near $r=1.1$, with
$\mathbf{r_{\max}=16}$ because beyond it every compressor's CTR is
near zero on family-a. \textbf{$50$ instances per family, $5$ seeds
per cell} ($250$ pairs per cell) is the smallest sample at which the
workload-level bootstrap CI clears the H1/H2 bands at the budgeted
compute; there is no formal a-priori power analysis, named as a
limitation in \charef{summary}. \textbf{$\theta_q$ at
$0.5\times p_0$} reads off ``coordination has collapsed'' without a
separate calibration sweep and sits below the H2 $30\%$-drop bar; the
CAAC ablation finds the coordination plateau invariant to
$\theta_q\in\{0.6,0.7,0.8\}$. The \textbf{bootstrap} uses BCa
intervals~\cite{efron1993bootstrap} at $n_\text{boot}=10{,}000$ on
the verdict pipelines and $500$ on the $\theta_q$ and
frontier-$\tau^{*}$ pipelines, whose per-iteration curve fit is the
budget driver.


\section{Scope of Evaluation: the Multi-Fragment Proxy}
\label{sec:scope}

The empirical protocol uses two planners, isolating the
compression-quality signal from LLM variance: a deterministic regex
information-extraction solver for H1/H2, and a single LLM call with
all compressed fragments visible for Corollary~1, the frontier
validation, and the real-data transfer (demonstrating that the cliff
structure the solver measures appears identically under a non-trivial
language model). The AutoGen v0.4 multi-round
backend~\cite{wu2023autogen} in \texttt{src/m6/orchestrator/}
integrates with the bus but is not used: multi-round negotiation
variance dominates the compression signal at the swept ratios, so the
backend demonstrates that the bus is multi-agent-ready rather than
serving as the measured planner. Coordination success
(Section~\ref{sec:c1}) is therefore the structural property that the
planner must combine information across fragments, not multi-round
communication quality, which \charef{summary} names as future work.

\section{Compute Envelope and Reproducibility}
\label{sec:envelope}

The empirical work runs on an Apple~M4~Pro 48~GB workstation (local
Ollama planners, deterministic-solver pipeline, and every figure when
the compression cache is in place) and an RTX~5090 32~GB host
(compression precompute, the $8$B model-scaling sweep, the frontier
API sweep, and the HotpotQA / MultiHop-RAG transfer). Total wallclock
for the canonical pipeline is $\sim 30$~GPU-hours plus $\sim 12$~hours
of laptop time. The reproducibility package is the complete
repository: a \texttt{docker-compose.yml} for the memory-bus
service, a release tag, model and data cards, and one-command
\texttt{make} targets for every headline figure; one CSV per
hypothesis run plus a \texttt{manifest.yaml} recording the resolved
configuration, the config hash, the git revision, and the platform.
Appendix~\ref{appendix} lists the commands and the full compute stack.
